\chapter{Comparing the Two Datasets: flexset and soloq}

The analysis is built on datasets collected from two different queue types: the different matchmaking mechanics of Flex Queue and Solo/Duo Queue produce inherently different graph structures. Running the same pipeline on both datasets means that any observed difference follows directly from the queue mechanics --- not from differences in model parameters.

\section{Structural Differences Between the Two Datasets}

Flex Queue (queue ID 440) allows team-up: players can form premade teams of 2 to 5. This means that some player pairs appearing repeatedly on the ally side represent \emph{intentional repetition}. The ally/enemy ratio (73.17\% ally in \textit{flexset}) is partly a consequence of this mechanism.

Solo/Duo Queue (queue ID 420) allows premades of at most two players. The ally/enemy ratio (65.24\% ally in \textit{soloq}) is lower, consistent with more randomly assigned matches producing fewer dominant ally connections.

\begin{table}[H]
\centering
\caption{Basic sizes of the two datasets}
\begin{tabularx}{\textwidth}{|>{\raggedright\arraybackslash}p{5cm}|r|r|}
\hline
\textbf{Metric} & \textbf{flexset} & \textbf{soloq} \\
\hline
Match JSON files & 4981 & 2001 \\
Raw players & 15\,147 & 2822 \\
Refined players & 15\,421 & 6768 \\
Clusters (DB) & 599 & 140 \\
Cluster members (DB) & 5741 & 1538 \\
Visible nodes & 5741 & 1538 \\
Visible edges & 18\,548 & 4111 \\
Average edge weight & 3.405 & 2.364 \\
Ally/enemy ratio & 73.17\% / 26.83\% & 65.24\% / 34.76\% \\
Visual clusters & 1531 & 870 \\
Average cluster size & 3.750 & 1.768 \\
Largest cluster size & 20 & 20 \\
\hline
\end{tabularx}
\end{table}

The average cluster size in \textit{flexset} (3.750) is more than twice that of \textit{soloq} (1.768): the former is dominated by groups of 3--4, the latter by duos of 1--2. The queue mechanics directly explain this: Flex Queue allows premades of up to 5, while Solo/Duo Queue allows a maximum of 2.

\section{Summary}

\begin{table}[H]
\centering
\caption{Summary comparison of the two datasets}
\begin{tabularx}{\textwidth}{|>{\raggedright\arraybackslash}p{5cm}|X|X|}
\hline
\textbf{Dimension} & \textbf{flexset} & \textbf{soloq} \\
\hline
Queue mechanics & Flex, premade-focused & Solo/Duo, individual-centred \\
Average edge weight & 3.405 (denser) & 2.364 (sparser) \\
Ally ratio & 73.17\% & 65.24\% \\
Average cluster size & 3.750 & 1.768 \\
Triangles/node\textsuperscript{†} & 4.95 & 1.03 \\
Core-periphery structure\textsuperscript{‡} & Bridge-rich associative core with many peripheral ally islands & Smaller, more randomly distributed local groups \\
Social-path opscore ass. & 0.2346 & 0.1750 \\
Battle-path opscore ass. & 0.1461 & 0.1625 \\
Feedscore battle-path ass. & $-0.0108$ & $-0.0366$ \\
Research role & Associative graph and performance relationship & Performance baseline, control \\
\hline
\end{tabularx}
\end{table}

\noindent\small{†~\textit{Triangle structures} (triads): if player A knows B and C, and B also knows C, this forms a closed triangle in the network. More such closed triangles indicate stronger group cohesion. \quad ‡~\textit{Core-periphery}: the inner core of the graph is densely connected, while peripheral nodes have only a few connections.}\normalsize

\medskip
\textit{flexset} is the more natural setting for the associative player graph and recurring ally connections; \textit{soloq} serves as the performance baseline and control dataset. The \textit{battle-path opscore} value is higher in \textit{soloq} (0.1625) than in \textit{flexset} (0.1461). In \textit{social-path} mode \textit{flexset} is stronger (0.2346 vs. 0.1750), reflecting the premade-ally mechanism.
